\uuid{usQQ}
\titre{Tangentes}
\auteur{hanani ; mzari}
\organisation{exo7}
\licence{CC BY-NC-SA 4.0}
\datecreate{2020-01-20}
\niveau{L1}
\module{Géométrie}
\chapitre{Courbes planes}
\sousChapitre{Courbes paramétrées}
\difficulte{1}
\theme{Courbes paramétrées, tangentes}
\exo7qcmid{lille-2.courbes.002}

\contenu{
\begin{question}
\qclassification{244.01}
\qtags{motcle=tangentes, points doubles}
La trajectoire $\Gamma$ d'une particule en mouvement est donnée par les équations
$$x=1+\cos (2t)\quad \mbox{et} \quad y=1+\sin (2t)\qquad 0\leq t\leq \pi.$$
\begin{answers}
\bad{$\Gamma$ est le cercle de centre $(1,1)$ et de rayon $2$.}
\good{$\Gamma$ possède un point double.}
\good{La tangente à $\Gamma$ au point $(x(0),y(0))$ est la droite d'équation $x=2$.}
\bad{La tangente à $\Gamma$ au point $(x(\pi),y(\pi))$ est la droite d'équation $y=x$.}
\end{answers}
\begin{explanations}
En éliminant le paramètre $t$, on obtient : $(x-1)^2+(y-1)^2=1$. Donc $\Gamma$ est le cercle de centre $(1,1)$ et de rayon $1$. Les points de paramètres $0$ et $\pi$ sont confondus. Enfin, $x'(0)=0$ et $y'(0)\neq 0$, donc la tangente au point $(x(0),y(0))$ est la droite d'équation $x=x(0)$.
\end{explanations}
\end{question}
}
